\paragraph{What the residual lines imply about the head.}
\hypertarget{source-lift-weight-exchange}{}
Let \(H=h_0\) and \(K=\lfloor H/(p-1)\rfloor\).
Here \(p>y>8\), \(p\) is unramified in the CM field, and
the large-prime inequalities give \(t+2K<p\).
Use all embeddings of the CM field at \(p\). At the conjugate embedding use the decreasing sequence obtained by reflection and reversal. In zero-based coordinates the operation is
\[
 (u^c)_a=H-u_{r-1-a}\quad(0\le a<r),\qquad (u^c)^c=u .
\] The
bottom of the new block is \(H-(B+t-1)\). Write
\(\sigma i\) for the source label of the Frobenius matrix
ending at \(i\).

In the local cyclic indexing it is \(i-1\).
This permutation commutes with conjugation. The Galois attachment can use either
orientation for the whole family. Pulling back the two
sequences and the block along that involution identifies its
\(J\) with \(U'\), and its line weights with
\(b_{i,s}=B_i+s\) at each embedding \(i\). A calculation at
the transported label then reads at \(i,\sigma i\) in the
original labels. Reflection preserves the characteristic-polynomial comparison. Indeed, for any list of residues of length \(r\),
\[
 \prod_{a=0}^{r-1}\bigl(X-(H-u_a)\bigr)=(-1)^r\prod_{a=0}^{r-1}\bigl((H-X)-u_a\bigr).
\]

\begin{proposition}[A depth bound from the residual lines]\label{prop:residual-head-depth}
For the shifted highest weights, conjugate labels and consecutive
line block just specified, the deficit at a head cut satisfies
\[
 \sum_{a<k_c+j}(\lambda_{i,a}-\mu_{i,a})
   \le qH+(p m_{\sigma i}+m_i)(t+2K).
\]
Here \(m_i\) counts the original shifted weights outside the consecutive
block at \(i\) whose residues occur in the block enlarged by \(K\)
at each end. The count \(m_{\sigma i}\) uses the Frobenius source label. At the indices satisfying the shape and
prime exclusions of Section~\ref{subsec:lift-exponent-primes},
both corrections vanish and the depth is at most \(\lfloor qh\rfloor\),
uniformly over the occurring Hecke ideal and the weight \(\mu\).
\end{proposition}

\begin{proof}
Use the norm coordinate of Section~\ref{subsec:cut-wach-models},
written \(X=Z\) in this calculation, with the same terminal
Hodge labels. Apply its construction to
the lattice of the congruent weight \(\mu\), to the inverse
images of the residual flag and to the specified scalar lines.
Their coefficient extension is over the selected place. Their
jumps are respectively \(U'_{i,a}\) and \(b_{i,s}\) and lie
in \([0,H]\). Write \(F\) for the residue field. In the resulting
\(F[[X]]\) lattice the reduced Frobenius has an integral right
inverse up to \(X^H\). For any \(d\)-rowed minor \(M\) of the reduced Frobenius matrix the elementary divisors give
\[
 \operatorname{ord}_X M\ \ge\ s_i(d),\qquad s_i(d)=\sum_{a=r-d}^{r-1}U'_{i,a}.
\]
The first cyclotomic matrix coefficient \(\Theta_i\) has
characteristic roots \(-\bar c U'_{i,a}\).
The residual filtration gives the same ranks and marked lines in
\(F((X))\) by exact 'etale realization.

Its intersection with
\(F[[X]]^r\) is a saturated flag, so extending bases at successive
steps gives one basis adapted to it. Each marked line occupies its own
diagonal block. Its two chosen generators in the Laurent
line differ by a unit times \(X^{\kappa_{i,s}}\), with the
sign fixed so that Frobenius in the adapted basis has order
\[
        \alpha_{i,s}=b_{i,s}+p\kappa_{\sigma i,s}-\kappa_{i,s}.
\]
In fact Frobenius sends the parameter to \(X^p\), and the
reduced rank-one Frobenius is a unit times \(X^{b_{i,s}}\).
The latter follows from the rank-one determinant, whose
order before reduction is \(b_{i,s}\) at \(X-p\) by height
and the local elementary divisor. Integrality gives
\(\alpha_{i,s}\ge0\) and the height inverse gives
\(\alpha_{i,s}\le H\) on the singleton block. The inverse
is block triangular after inverting \(X\), so the two
diagonal entries multiply to \(X^H\).

\proofstep{From the determinant to the minor estimate.}
To constrain individual head weights we pass from the determinant to a minor.

For the cyclotomic generator, \(\gamma(X)/X\) has first
terms \(1+\bar c X\) modulo \(X^2\) in this reduced
coordinate. For a unit \(v=v_0+v_1X+O(X^2)\), the constant and linear terms cancel in the quotient:
\[
 \gamma(v)\equiv v_0+v_1X\equiv v\pmod{X^2},\qquad \gamma(v)/v\equiv1\pmod{X^2}.
\] The coefficient of \(X\) on the same marked line
is therefore \(-\bar c(b_{i,s}-\kappa_{i,s})\).
An integral basis change conjugates this first coefficient by
its constant matrix, since the cyclotomic action fixes the
change of basis modulo \(X^2\).

Taking characteristic
polynomials of the block triangular matrix and cancelling
the nonzero common scalar shows that the residues
\(b_{i,s}-\kappa_{i,s}\), counted with multiplicity, occur
among those of \(U'_i\). The determinant of the Frobenius
minor on the marked positions is the product of their
diagonal entries. The minor divisibility then gives
\begin{equation}\label{eq:lift-line-minor}
       s_i(t)\le\sum_s b_{i,s}
              +p\sum_s\kappa_{\sigma i,s}-\sum_s\kappa_{i,s}.
\end{equation}
The flag and the rank-one comparison here come from the
same residual representation at the same maximal ideal.
\label{cor:lift-weight-exchange-1}

Since \(b_{i,s}\) and \(\alpha_{i,s}\) belong to \([0,H]\),
we have \(|p\kappa_{\sigma i,s}-\kappa_{i,s}|\le H\).
For fixed \(s\), let \(j\) maximize \(\lvert\kappa_{j,s}\rvert\).
Apply the relation at \(i=\sigma^{-1}j\), where the term
with coefficient \(p\) is \(\kappa_{j,s}\). The triangle inequality at this predecessor gives
\[
 p|\kappa_{j,s}|\le H+|\kappa_{\sigma^{-1}j,s}|
            \le H+|\kappa_{j,s}| .
\]
Hence every \(|\kappa_{i,s}|\le H/(p-1)\).
To bound the sum of the shifts, enlarge the block
\(I=[B_i,B_i+t-1]\) by \(K\) at both ends to \(I^+\).

Let \(m_i\) count original weights outside \(I\) whose
residue occurs in \(I^+\). The strict length bound makes
reduction modulo \(p\) injective on the integers in \(I^+\).
The \(t\) shifted entries \(b_{i,s}-\kappa_{i,s}\) lie
there. Match their residues to distinct occurrences among
the \(U_i\) using the two multiset comparisons. At most
\(m_i\) are matched outside \(I\). All other matches are
equalities in \(I^+\). Pair the remaining shifted entries
arbitrarily with the unused block entries. Each difference
is bounded by \(t+2K\), so
\[
                \left|\sum_s\kappa_{i,s}\right|
                            \le m_i(t+2K).
\]
This also holds at \(\sigma i\), where the exceptional count
uses its own block. Reflection preserves the count because it
carries the block, its enlargement and residue congruence
bijectively to their reflected versions.

\proofstep{Returning to the head of the weight.}
The deficit of the first \(k_c+j\) weights equals the excess
of the rest by equality of total sums. Of the last \(t+q\)
entries of \(U'_i\), the smallest \(t\) have sum \(s_i(t)\)
and the other \(q\) are at most \(H\) apiece. In \(U_i\)
the corresponding total is at least \(\sum_s b_{i,s}\)
because the bottom \(q\) entries are nonnegative. Equation~\eqref{eq:lift-line-minor}
thus proves
\[
 \sum_{a<k_c+j}(\lambda_{i,a}-\mu_{i,a})
   \le qH+(p m_{\sigma i}+m_i)(t+2K).
\]
The staircase cancels term by term on the left. The shape
and prime exclusions of Section~\ref{subsec:lift-exponent-primes}
give \(m_i=m_{\sigma i}=0\). Every surviving highest
weight at this index with a larger weight therefore has depth at most
\(i_i=\lfloor qh\rfloor\), with the choice of
\(\mu\) and of the occurring Hecke ideal unrestricted.
\label{cor:lift-weight-exchange-2}
\end{proof}
